\begin{pastexam}[label={px:ncku-111-prog-08}]{成大資訊聯招 111 學年度 第 8 題}
（填表）

\textbf{[4\%]} The following is a hash table with linear probing. Assume the hash function is f(x) = x mod 17 for indexing the hash table. Data items 23, 52, 11, 1, 50, 99, 65, 20, 35, 34 and 82 are inserted into the hash table in order. Please fill the five blanks in the table.
\begin{center}
\begin{tabular}{|l|l|}
\hline
key & value \\ \hline
0 & 34 \\ \hline
1 & (A) ? \\ \hline
2 & (B) ? \\ \hline
3 & (C) ? \\ \hline
4 & (D) ? \\ \hline
5 & (E) ? \\ \hline
6 & 23 \\ \hline
7 & \\ \hline
8 & \\ \hline
9 & \\ \hline
10 & \\ \hline
11 & 11 \\ \hline
12 & \\ \hline
13 & \\ \hline
14 & 99 \\ \hline
15 & 65 \\ \hline
16 & 50 \\ \hline
\end{tabular}
\end{center}
\end{pastexam}

\begin{pastsoln}
答案：(A) $52$，(B) $1$，(C) $20$，(D) $35$，(E) $82$。

線性探測（資結書第~12.3~節）：算出 $f(x) = x \bmod 17$，那一格有東西就往下一格找（$16$ 的下一格是 $0$），
放進第一個空格。 依序插入：
\begin{center}
\begin{tabular}{rrp{7.2cm}l}
\toprule
$x$ & $x \bmod 17$ & 探測過的格子 & 放入 \\
\midrule
23 & 6 & 6 & 6 \\
52 & 1 & 1 & 1 \\
11 & 11 & 11 & 11 \\
1 & 1 & 1（52） & 2 \\
50 & 16 & 16 & 16 \\
99 & 14 & 14 & 14 \\
65 & 14 & 14（99） & 15 \\
20 & 3 & 3 & 3 \\
35 & 1 & 1（52）、2（1）、3（20） & 4 \\
34 & 0 & 0 & 0 \\
82 & 14 & 14（99）、15（65）、16（50）、0（34）、1（52）、2（1）、3（20）、4（35） & 5 \\
\bottomrule
\end{tabular}
\end{center}
餘數的算法：$52 = 3 \times 17 + 1$，$50 = 2 \times 17 + 16$，$99 = 5 \times 17 + 14$，$65 = 3 \times 17 + 14$，
$35 = 2 \times 17 + 1$，$34 = 2 \times 17$，$82 = 4 \times 17 + 14$。

最後的表：第 $0$ 格 $34$、第 $1$ 格 $52$、第 $2$ 格 $1$、第 $3$ 格 $20$、第 $4$ 格 $35$、第 $5$ 格 $82$、
第 $6$ 格 $23$、第 $11$ 格 $11$、第 $14$--$16$ 格 $99$、$65$、$50$，和題目已經填好的格子一致。 $82$ 從
第 $14$ 格一路繞過表尾探測了 $9$ 次才找到空格，這就是線性探測的\textbf{主要群聚}：第 $14$--$16$ 格和
第 $0$--$6$ 格連成一大段，落在這段裡的鍵都要走到段尾。
\end{pastsoln}
